\chapter{Black--Scholes de trois façons}\label{ch:dv:black-scholes-three-ways}

Le 26 avril 1973, le Chicago Board Options Exchange ouvre ses portes avec
des options d'achat sur 16 actions et échange 911 contrats. Au même printemps, le
\emph{Journal of Political Economy} publie un article de Fischer Black et Myron
Scholes, accompagné d'un commentaire de Robert Merton, qui donne le prix d'une telle option
d'achat sous la forme d'une formule fermée à cinq paramètres~: le cours au comptant,
le prix d'exercice, la durée, le taux et la volatilité. La formule est juste pour une
raison qui surprend les économistes~: le rendement espéré de l'action, le seul chiffre
sur lequel tout investisseur débat, ne figure pas parmi les cinq. Ce chapitre établit
la formule de trois manières, car chaque démonstration apporte un enseignement différent~:
la réplication montre pourquoi la dérive disparaît et ce que fait réellement un hedger~;
la voie martingale montre quelle espérance est calculée et comment la généraliser~;
la limite binomiale montre que rien ici ne requiert le calcul stochastique en temps continu
pour être vrai, seulement pour être commode.

\section{Le modèle et l'argument de réplication}

\begin{definition}[Modèle de Black--Scholes]\label{def:dv:black-scholes-three-ways:model}
Le \emph{modèle de Black--Scholes}\index{modèle de Black--Scholes} est un marché
comportant un compte bancaire dont le montant croît au taux constant $r$ et une action qui
verse un rendement de dividende continu $q$ et suit un mouvement brownien géométrique (livre~4,
chapitre~4) sous la probabilité historique,
\[
dS_t=\mu S_t\,dt+\sigma S_t\,dW_t,
\]
avec $\mu$ constant et $\sigma>0$. Les transactions sont continues et sans frictions,
les ventes à découvert sont autorisées, et les liquidités peuvent être empruntées et prêtées au taux $r$.
\end{definition}

Soit $V(t,S)$ la valeur d'un produit européen versant $g(S_T)$, supposée régulière en
$(t,S)$. On détient le produit et l'on vend $\Delta$ actions. La position courte verse le
dividende $qS\,dt$ au prêteur de l'action. Par la formule d'Itô (livre~4,
chapitre~3), le portefeuille $\Pi=V-\Delta S$ varie de
\[
d\Pi=\Bigl(\partial_tV+\tfrac12\sigma^2S^2\partial_{SS}V-q\Delta S\Bigr)dt
+\bigl(\partial_SV-\Delta\bigr)\,dS .
\]
Le choix $\Delta=\partial_SV$ élimine l'unique terme aléatoire. Un portefeuille
sans risque doit rapporter le taux sans risque, sans quoi emprunter pour l'acquérir
(ou le vendre pour prêter) constituerait un arbitrage~: $d\Pi=r\Pi\,dt=r(V-S\partial_SV)\,dt$. L'égalité des
deux expressions de $d\Pi$ donne l'équation de valorisation.

\begin{definition}[Équation de Black--Scholes]\label{def:dv:black-scholes-three-ways:pde}
L'\emph{équation de Black--Scholes}\index{equation de Black--Scholes@équation de Black--Scholes} est une
équation aux dérivées partielles parabolique
\[
\partial_tV+(r-q)S\,\partial_SV+\tfrac12\sigma^2S^2\,\partial_{SS}V-rV=0,
\qquad V(T,S)=g(S),
\]
satisfaite par la valeur de tout produit européen sur l'action dans le
\omterm{def:dv:black-scholes-three-ways:model}{modèle de Black--Scholes}.
\end{definition}

Le terme de dérive $\mu$ multiplie $dS$, et $dS$ est couvert~: il s'annule avec la
position en actions. Ce qui subsiste, c'est le terme de variance $\tfrac12\sigma^2S^2
\partial_{SS}V$, qu'aucune position en actions ne peut annuler car il provient de la
courbure de $V$, ainsi que le financement de la couverture au taux $r-q$. L'équation est
linéaire, si bien que la valeur d'un portefeuille de produits est la somme des valeurs, et
elle est la même pour tout investisseur, quelle que soit son opinion sur $\mu$.

\begin{remark}[Ce que fait le hedger]\label{rem:dv:black-scholes-three-ways:hedger}
Vendre le produit et détenir $\Delta=\partial_SV$ actions financées au taux $r$ constitue une
\omterm{def:dv:no-arbitrage-and-the-fundamental-theorems:sf}{stratégie autofinancée} (chapitre~1) dont la valeur suit exactement $V$ dans le modèle.
La \Cref{fig:dv:black-scholes-three-ways:replication} la met en œuvre une fois, rééquilibrée
chaque jour le long d'une trajectoire simulée avec une dérive de 12~\% par an~: le portefeuille suit
l'option, et l'écart à l'échéance, 0.39 sur une prime de 10.45, est le coût
d'une couverture quotidienne plutôt que continue, objet du chapitre~4.
\end{remark}

\begin{omfigure}
\begin{tikzpicture}
\begin{axis}[width=11cm,height=5cm,
  xlabel={temps (années)},ylabel={valeur},
  tick label style={font=\small},label style={font=\small},
  xmin=0,xmax=1,ymin=0,ymax=26,grid=major,grid style={gray!25},
  legend columns=2,legend style={font=\footnotesize,at={(0.5,-0.3)},anchor=north,draw=none,fill=none,/tikz/every even column/.append style={column sep=8pt}},
  table/col sep=comma]
\addplot[omDef,very thick] table[x=t,y=option]{figdata/derivatives/03-black-scholes-three-ways/replication.csv};
\addplot[omThm,thick,dashed] table[x=t,y=portfolio]{figdata/derivatives/03-black-scholes-three-ways/replication.csv};
\legend{valeur de l'option par la formule,portefeuille de réplication (rééquilibrage quotidien)}
\end{axis}
\end{tikzpicture}

\omcaption{Réplication le long d'une trajectoire. Une option d'achat à la monnaie d'un an ($r=5\%$,
$\sigma=20\%$) est répliquée en détenant $\Delta$ actions et en empruntant, rééquilibrée
252 fois~; l'action, simulée avec une dérive de 12~\%, termine à 115.46. Les deux
courbes sont indiscernables à cette échelle~; l'écart final est de 0.39. Données~: le tutoriel.}\label{fig:dv:black-scholes-three-ways:replication}
\end{omfigure}

\section{L'équation de valorisation}

L'\omterm{def:dv:black-scholes-three-ways:pde}{équation de Black--Scholes} est une équation de la chaleur rétrograde déguisée. Avec
$x=\ln S$, la durée restante $\tau=T-t$ et $V=e^{-r\tau}u$, elle devient
$\partial_\tau u=(r-q-\tfrac12\sigma^2)\partial_xu+\tfrac12\sigma^2\partial_{xx}u$, une
diffusion à coefficients constants. Par la formule de Feynman--Kac (livre~4,
chapitre~4), sa solution est une espérance~:
\begin{equation}\label{eq:dv:black-scholes-three-ways:fk}
V(t,S)=e^{-r(T-t)}\,\E\bigl[g(S_T)\mid S_t=S\bigr],\qquad
dS_u=(r-q)S_u\,du+\sigma S_u\,d\widetilde W_u .
\end{equation}
L'espérance est prise pour une action qui croît au taux $r-q$, et non au taux $\mu$~:
l'équation ignore tout de $\mu$, et sa solution aussi. C'est d'où vient
l'expression «~risque-neutre~»~: les prix sont des espérances comme si les investisseurs
n'exigeaient aucune prime pour le risque. Il s'agit d'une affirmation sur les prix,
comme le chapitre~1 la lisait, et non sur les investisseurs.

\begin{omfigure}
\begin{tikzpicture}
\begin{axis}[width=10.5cm,height=5cm,
  xlabel={cours au comptant},ylabel={valeur de l'option},
  tick label style={font=\small},label style={font=\small},
  xmin=60,xmax=140,ymin=0,ymax=45,grid=major,grid style={gray!25},
  legend columns=4,legend style={font=\footnotesize,at={(0.5,-0.3)},anchor=north,draw=none,fill=none,/tikz/every even column/.append style={column sep=8pt}},
  table/col sep=comma]
\addplot[gray,thick,dashed] table[x=spot,y=payoff]{figdata/derivatives/03-black-scholes-three-ways/call_vs_spot.csv};
\addplot[omProp,thick] table[x=spot,y=t01]{figdata/derivatives/03-black-scholes-three-ways/call_vs_spot.csv};
\addplot[omDef,thick] table[x=spot,y=t05]{figdata/derivatives/03-black-scholes-three-ways/call_vs_spot.csv};
\addplot[omThm,very thick] table[x=spot,y=t1]{figdata/derivatives/03-black-scholes-three-ways/call_vs_spot.csv};
\legend{gain à l'échéance,0.1 an,0.5 an,1 an}
\end{axis}
\end{tikzpicture}

\omcaption{Le call de strike 100 comme solution de l'équation de pricing, à
trois échéances ($r=5\%$, $\sigma=20\%$). En remontant le temps, le coude du payoff
au strike est lissé par la diffusion et la courbe est relevée par les
intérêts sur le strike. Données~: le code du chapitre.}\label{fig:dv:black-scholes-three-ways:values}
\end{omfigure}

\section{La voie martingale}

Le chapitre~1 évaluait un \omterm{def:dv:no-arbitrage-and-the-fundamental-theorems:claim}{actif contingent} comme une espérance actualisée sous une mesure dans laquelle
les prix actualisés des actifs négociés sont des martingales. Dans le \omterm{def:dv:black-scholes-three-ways:model}{modèle de Black--Scholes}, le
théorème de Girsanov (livre~4, chapitre~5) le donne explicitement~: sous la
probabilité risque-neutre $\mathbb Q$, le processus $W^{\mathbb Q}_t=W_t+\frac{\mu-r+q}
{\sigma}t$ est un mouvement brownien et
\[
S_T=S_0\exp\Bigl(\bigl(r-q-\tfrac12\sigma^2\bigr)T+\sigma W^{\mathbb Q}_T\Bigr),
\]
de sorte que $e^{-(r-q)t}S_t$ est une $\mathbb Q$-martingale~: l'action ajustée des dividendes,
actualisée au taux $r$, est un jeu équitable. Le prix du call est
\[
C=e^{-rT}\E^{\mathbb Q}\bigl[(S_T-K)\mathbf 1_{S_T>K}\bigr]
=e^{-rT}\E^{\mathbb Q}\bigl[S_T\mathbf 1_{S_T>K}\bigr]-Ke^{-rT}\,\mathbb Q(S_T>K).
\]
Le second terme est une probabilité normale. Le premier demande une idée de plus~: changer
de numéraire pour l'action (avec réinvestissement de ses dividendes). Sous la
mesure de l'action $\mathbb Q^S$, de densité $S_Te^{qT}/(S_0e^{rT})$, la première espérance
devient $S_0e^{(r-q)T}\,\mathbb Q^S(S_T>K)$, et sous $\mathbb Q^S$ le logarithme de l'action
a pour dérive $r-q+\tfrac12\sigma^2$ au lieu de $r-q-\tfrac12\sigma^2$.

\begin{theorem}[Formule de Black--Scholes]\label{thm:dv:black-scholes-three-ways:formula}
Dans le \omterm{def:dv:black-scholes-three-ways:model}{modèle de Black--Scholes}, le call et le put européens valent
\[
C=Se^{-qT}\Phi(d_1)-Ke^{-rT}\Phi(d_2),\qquad
P=Ke^{-rT}\Phi(-d_2)-Se^{-qT}\Phi(-d_1),
\]
\[
d_{1,2}=\frac{\ln(S/K)+(r-q)T}{\sigma\sqrt T}\pm\tfrac12\sigma\sqrt T,
\]
la \emph{formule de Black--Scholes}\index{formule de Black--Scholes}. Ici
$\Phi(d_2)=\mathbb Q(S_T>K)$ est la probabilité risque-neutre d'exercice et
$\Phi(d_1)=\mathbb Q^S(S_T>K)$ la même probabilité sous la mesure de l'action, et
$e^{-qT}\Phi(d_1)=\partial_SC$ est le ratio de couverture.
\end{theorem}
\begin{proof}
$\ln S_T$ est normal de variance $\sigma^2T$ sous les deux mesures, avec des moyennes
qui diffèrent de $\sigma^2T$~; $\mathbb Q(S_T>K)=\Phi(d_2)$ et $\mathbb Q^S(S_T>K)=
\Phi(d_1)$ s'obtiennent en centrant et réduisant. Le put se déduit de la parité call--put. Enfin,
l'identité $\partial_SC=e^{-qT}\Phi(d_1)$ utilise $S e^{-qT}\varphi(d_1)=Ke^{-rT}
\varphi(d_2)$, qui annule les termes issus de la dérivation de $d_1$ et $d_2$.
\end{proof}

\begin{example}[Le call de référence du desk]\label{ex:dv:black-scholes-three-ways:ref}
$S=K=100$, $T=1$, $r=5\%$, $q=0$, $\sigma=20\%$~: $d_1=0.35$, $d_2=0.15$,
$\Phi(d_1)=0.6368$, $\Phi(d_2)=0.5596$, $C=10.4506$, $P=5.5735$. La probabilité risque-neutre
d'exercice vaut 0.56~; avec une dérive réelle de 12~\% la
probabilité que le call finisse dans la monnaie vaut 0.69, et elle ne joue aucun rôle dans le
prix.
\end{example}

\begin{proposition}[Approximation à la monnaie]\label{prop:dv:black-scholes-three-ways:atm}
Au forward, $K=F=Se^{(r-q)T}$, le call et le put valent tous deux
$e^{-rT}F\bigl(2\Phi(\tfrac12\sigma\sqrt T)-1\bigr)\approx0.4\,e^{-rT}F\sigma\sqrt T$.
\end{proposition}
\begin{proof}
Avec $K=F$, $d_{1,2}=\pm\tfrac12\sigma\sqrt T$ et $C=e^{-rT}F(\Phi(d_1)-\Phi(-d_1))$~;
on développe $\Phi$ en 0 avec la pente $\varphi(0)=1/\sqrt{2\pi}\approx0.3989$.
\end{proof}

\section{La limite binomiale}

L'\omterm{def:dv:the-binomial-model:crr}{arbre de Cox--Ross--Rubinstein} du chapitre~2 ne demande aucun calcul stochastique et
atteint le même nombre. Avec $n$ pas, le call sur arbre est l'espérance actualisée de
$(S_T-K)^+$ sous une loi binomiale du nombre de mouvements à la hausse. Quand
$n\to\infty$ avec $u=e^{\sigma\sqrt{T/n}}$, la probabilité de l'arbre de finir au-dessus de
$K$ converge vers $\Phi(d_2)$, et la même probabilité calculée avec les poids
$pu/R$ et $(1-p)d/R$ (la mesure de l'action sur l'arbre) converge vers
$\Phi(d_1)$~: les deux termes de la formule apparaissent comme les limites des deux sommes binomiales de la
\cref{prop:dv:the-binomial-model:formula}. Pour le call de référence, l'arbre à
1\,000 pas donne 10.4486, soit 0.0020 en dessous de la formule, l'erreur étant d'ordre
$1/n$ comme mesuré au chapitre~2.

\section{La formule de Black, et ce qu'enseigne chaque dérivation}

De nombreux sous-jacents sont cotés en termes de forwards ou de contrats à terme~: contrats à terme sur matières premières,
contrats à terme obligataires, taux forwards, forwards de change. Évaluer sur le forward retire de la formule le
rendement de dividende, le coût d'emprunt et le financement~; tous sont
inclus dans $F$.

\begin{definition}[Modèle de Black]\label{def:dv:black-scholes-three-ways:black}
Le \emph{modèle de Black}\index{modèle de Black} suppose que le prix forward (ou du contrat à terme)
$F_t$ du sous-jacent pour livraison en $T$ est un mouvement brownien géométrique
sans dérive sous la $T$-probabilité forward (livre~4, chapitre~5)~:
$dF_t=\sigma F_t\,dW^T_t$. Un call européen d'échéance $T$ vaut alors
\[
C=P(0,T)\bigl(F\Phi(d_1)-K\Phi(d_2)\bigr),\qquad d_{1,2}=\frac{\ln(F/K)}{\sigma\sqrt T}\pm\tfrac12\sigma\sqrt T .
\]
\end{definition}

La \omterm{thm:dv:black-scholes-three-ways:formula}{formule de Black--Scholes} est le cas particulier $F=Se^{(r-q)T}$,
$P(0,T)=e^{-rT}$. Black a publié la version forward en 1976 pour les
contrats sur matières premières~; les options sur contrats à terme du livre~3, chapitre~12, les caplets
du livre~2, chapitre~13, et les options de change du chapitre~20 s'évaluent toutes
avec elle, et la volatilité implicite se définit à travers elle (livre~1,
chapitre~25). Quand le taux est stochastique, actualiser avec $P(0,T)$ et prendre
l'espérance sous la probabilité forward conserve la formule exacte, à condition que
le forward lui-même soit log-normal.

\begin{center}
\footnotesize
\begin{tabular}{@{}p{2.4cm}p{4.9cm}p{5.4cm}@{}}
\toprule
Voie & Ce qu'elle suppose & Ce qu'elle enseigne \\
\midrule
Réplication & formule d'Itô, couverture continue & pourquoi $\mu$ disparaît~; le ratio de couverture~; l'équation se généralise à tout payoff et aux grilles numériques (chapitre~22) \\
Martingale & une mesure, un numéraire & quelle espérance est prise~; le choix du numéraire simplifie les payoffs~; se généralise à Monte Carlo (chapitre~23) et à d'autres dynamiques (chapitres~9 à 13) \\
Limite binomiale & deux états par pas & que le résultat relève de l'arbitrage, non du calcul~; l'exercice américain (chapitre~6) \\
\bottomrule
\end{tabular}
\end{center}

\begin{omfigure}
\begin{tikzpicture}[font=\small,
  bx/.style={draw=omDef,thick,rounded corners=2pt,fill=omDef!6,align=center,text width=3.3cm,inner sep=4pt},
  res/.style={draw=omThm,very thick,rounded corners=2pt,fill=omThm!8,align=center,text width=3.6cm,inner sep=5pt}]
\node[bx] (a) at (0,2.2) {détenir $\partial_SV$ actions~:\\Itô élimine $dW$};
\node[bx] (b) at (0,0) {Girsanov~: $\mathbb Q$, puis\\le numéraire action};
\node[bx] (c) at (0,-2.2) {arbre à $n$ pas,\\théorème central limite};
\node[bx] (a2) at (4.6,2.2) {équation de prix\\+ Feynman--Kac};
\node[bx] (b2) at (4.6,0) {$\mathbb Q(S_T>K)=\Phi(d_2)$\\$\mathbb Q^S(S_T>K)=\Phi(d_1)$};
\node[bx] (c2) at (4.6,-2.2) {sommes binomiales\\$\to$ intégrales gaussiennes};
\node[res] (f) at (9.3,0) {$Se^{-qT}\Phi(d_1)$\\$-\,Ke^{-rT}\Phi(d_2)$};
\draw[-{Stealth},thick] (a) -- (a2);
\draw[-{Stealth},thick] (b) -- (b2);
\draw[-{Stealth},thick] (c) -- (c2);
\draw[-{Stealth},thick] (a2.east) -- (f.north west);
\draw[-{Stealth},thick] (b2.east) -- (f.west);
\draw[-{Stealth},thick] (c2.east) -- (f.south west);
\end{tikzpicture}

\omcaption{Trois voies vers une même formule~: la réplication (en haut), l'espérance sous la martingale
(au milieu), la limite de l'arbre binomial (en bas). Chaque voie mobilise un outil
mathématique différent, et chacune se généralise à une méthode
numérique différente.}\label{fig:dv:black-scholes-three-ways:routes}
\end{omfigure}

\begin{method}[Volatilité implicite, de façon robuste]\label{met:dv:black-scholes-three-ways:iv}
Pour inverser la formule de Black à partir d'un prix $c$~:
\begin{enumerate}
\item Vérifier $P(0,T)(F-K)^+\le c<P(0,T)F$ (call)~; hors de ces bornes, il n'existe
aucune volatilité, et la cotation ou le forward est erroné.
\item Encadrer~: le prix est croissant en $\sigma$~; doubler une borne supérieure jusqu'à ce qu'elle
fournisse un prix supérieur à $c$.
\item Partir d'une estimation en forme fermée sur la valeur temps, puis exécuter la méthode de Newton
avec le vega comme dérivée, en revenant à la dichotomie chaque fois qu'un pas sort de
l'encadrement. Arrêter lorsque l'erreur relative sur le prix atteint $10^{-12}$.
\end{enumerate}
Newton seul échoue loin de la monnaie, où le vega est très faible et un pas dépasse la cible~;
les bibliothèques de production recourent à des approximations rationnelles qui convergent en deux
itérations jusqu'à la précision machine (Jäckel).
\end{method}

\section{Tutoriel~: un prix, trois voies}\label{tut:dv:black-scholes-three-ways:tut}

\begin{tutorial}
\textbf{Objectif.} Implémenter la formule de Black et un solveur robuste de volatilité implicite,
puis pricer le call de référence par la formule, par un arbre et par Monte Carlo, et le répliquer
le long d'une trajectoire. \textbf{État final~:}
\cref{fig:dv:black-scholes-three-ways:replication,fig:dv:black-scholes-three-ways:mc}
et les nombres du problème du week-end.
\begin{tutsteps}
\item \textbf{La formule}, sur le forward, avec une fonction de répartition normale précise dans les queues~:
\omcode{code/firm/bs/firm_bs.py}{23}{36}{Formule de Black sur un forward, et Black--Scholes sur un spot.}\label{lst:dv:black-scholes-three-ways:black}
\item \textbf{Le solveur} de la \cref{met:dv:black-scholes-three-ways:iv}~:
\omcode{code/firm/bs/firm_bs.py}{61}{94}{Volatilité implicite~: bornes, encadrement, estimation initiale, Newton protégé.}\label{lst:dv:black-scholes-three-ways:iv}
\item \textbf{Monte Carlo}~: tirer $S_T$ exactement selon sa loi lognormale, moyenner le
payoff actualisé, et indiquer l'erreur type.
\omcode{code/derivatives/03-black-scholes-three-ways/python/dv_blackscholes.py}{24}{34}{Monte Carlo sous la probabilité risque-neutre, avec son erreur type.}\label{lst:dv:black-scholes-three-ways:mc}
\item \textbf{Exécuter} \texttt{dv\_blackscholes.problem()} et
\texttt{fig\_blackscholes.py}~; les jumeaux C++20 et Rust dans
\texttt{code/firm/bs/} passent les mêmes tests d'aller-retour.
\end{tutsteps}
\textbf{Pour aller plus loin.} Activer les variables antithétiques et comparer l'erreur type
à effort de calcul égal~; rééquilibrer le \omterm{def:dv:no-arbitrage-and-the-fundamental-theorems:claim}{portefeuille de réplication} chaque semaine et chaque heure, et
comparer l'écart à l'échéance avec celui obtenu au quotidien.
\end{tutorial}

\begin{omfigure}
\begin{tikzpicture}
\begin{semilogxaxis}[width=10.5cm,height=5cm,
  xlabel={nombre de trajectoires $M$ (échelle log)},ylabel={prix du call},
  tick label style={font=\small},label style={font=\small},
  xmin=80,xmax=1.3e6,ymin=7.5,ymax=15.5,grid=major,grid style={gray!25},
  legend columns=3,legend style={font=\footnotesize,at={(0.5,-0.3)},anchor=north,draw=none,fill=none,/tikz/every even column/.append style={column sep=8pt}},
  table/col sep=comma]
\addplot[name path=hi,omDef!40,thin] table[x=m,y=hi]{figdata/derivatives/03-black-scholes-three-ways/mc_convergence.csv};
\addplot[name path=lo,omDef!40,thin] table[x=m,y=lo]{figdata/derivatives/03-black-scholes-three-ways/mc_convergence.csv};
\addplot[omDef!20] fill between[of=hi and lo];
\addplot[omDef,very thick,mark=*,mark size=1.5pt] table[x=m,y=estimate]{figdata/derivatives/03-black-scholes-three-ways/mc_convergence.csv};
\addplot[omThm,thick,dashed] table[x=m,y=formula]{figdata/derivatives/03-black-scholes-three-ways/mc_convergence.csv};
\legend{,,deux erreurs types,estimation Monte Carlo,formule 10.4506}
\end{semilogxaxis}
\end{tikzpicture}

\omcaption{Estimations Monte Carlo du call de référence sur les $M$ premiers tirages parmi
un million. La bande de deux erreurs types se resserre en $1/\sqrt M$~: cent fois plus de
trajectoires rapportent une décimale significative de plus. Données~: le
tutoriel.}\label{fig:dv:black-scholes-three-ways:mc}
\end{omfigure}

\section{Construction~: les noyaux de Black--Scholes}\label{bld:dv:black-scholes-three-ways:bs}

\begin{build}
\bfield{Rôle} La boucle interne de la firme miniature~: chaque ajustement de surface, chaque
rapport de grecques, le moteur de risque du livre~6 et la bibliothèque de pricing
du chapitre~28 appellent ces fonctions des millions de fois par jour, si bien qu'elles existent
en Python, en C++20 et en Rust avec les mêmes tests.
\bfield{Interface} \texttt{black(fwd, strike, t, df, vol, right)}~;
\texttt{bs(spot, strike, t, r, q, vol, right)}~; \texttt{greeks(\ldots)} renvoyant
delta, gamma, vega, theta, rho, vanna et volga~;
\texttt{implied\_vol(price, fwd, strike, t, df, right)}. C++~:
\texttt{firm::bs} dans \texttt{cpp/firm\_bs.hpp}~; Rust~: crate \texttt{firm\_bs}.
\bfield{Règles} Pricer sur le forward~; $\Phi$ calculée par la fonction d'erreur complémentaire, afin que les prix
profondément hors de la monnaie conservent leur précision relative~; un prix hors des
bornes d'absence d'arbitrage lève une exception et ne renvoie jamais une volatilité.
\bfield{Tests d'acceptation} \texttt{code/firm/bs/tests/} et les tests C++ et Rust~:
la valeur du manuel 10.450583572185579~; la parité avec un rendement de dividende~; chaque
grecque comparée à une différence centrée~; des allers-retours de volatilité implicite à une erreur relative de
$10^{-11}$ sur cinq décennies de moneyness et des maturités allant d'un jour à 30 ans.
\bfield{Pour aller plus loin} Remplacer le solveur par une méthode d'approximation rationnelle et
mesurer le nombre d'itérations et le temps par inversion dans les trois langages.
\end{build}

\begin{omsources}
\item F.~Black and M.~Scholes, «~The pricing of options and corporate
liabilities~», \emph{Journal of Political Economy} 81 (1973) 637--654.
\item R.~C.~Merton, «~Theory of rational option pricing~», \emph{Bell Journal of
Economics and Management Science} 4 (1973) 141--183.
\item F.~Black, «~The pricing of commodity contracts~», \emph{Journal of
Financial Economics} 3 (1976) 167--179.
\item P.~Jäckel, «~Let's be rational~», \emph{Wilmott} (2015) 40--53.
\item H.~A.~Baker, «~Cboe at 50~», \emph{Financial History} (printemps 2023).
\end{omsources}

\section{Exercices}

\begin{exercise}[$\star$]\label{exo:dv:black-scholes-three-ways:1}
Calculer le call et le put à six mois, de prix d'exercice 100, sur une action valant 100 avec $r=3\%$,
sans dividende et avec $\sigma=25\%$.
\end{exercise}

\begin{exercise}[$\star$]\label{exo:dv:black-scholes-three-ways:2}
Vérifier que $V=S$ (en détenant l'action, avec $q=0$) et $V=e^{rt}$ (le compte bancaire)
résolvent l'\omterm{def:dv:black-scholes-three-ways:pde}{équation de Black--Scholes}. Que représente $V=Se^{-q(T-t)}$ lorsque
$q>0$~?
\end{exercise}

\begin{exercise}[$\star$]\label{exo:dv:black-scholes-three-ways:3}
Un contrat à terme se négocie à 80. Calculer un call à trois mois, de prix d'exercice 85, avec un
facteur d'actualisation de 0.99 et $\sigma=35\%$.
\end{exercise}

\begin{exercise}[$\star\star$]\label{exo:dv:black-scholes-three-ways:4}
Un trader anticipe que l'action sous-jacente au call de référence rapportera 15~\% par an et
en conclut que le call est sous-évalué à 10.4506. Que devrait croire le trader pour que le
call soit mal évalué, et comment le trader en tirerait-il profit~?
\end{exercise}

\begin{exercise}[$\star\star$]\label{exo:dv:black-scholes-three-ways:5}
Pour le call de référence, donner la probabilité risque-neutre d'exercice et la
probabilité réelle avec $\mu=12\%$. Laquelle des deux le prix utilise-t-il, et pourquoi
aucune des deux n'est-elle le ratio de couverture~?
\end{exercise}

\begin{exercise}[$\star\star$]\label{exo:dv:black-scholes-three-ways:6}
Dériver le call par rapport au prix d'exercice pour trouver le prix d'un
call digital qui verse 1 si $S_T>K$, et l'évaluer pour les paramètres de référence.
\end{exercise}

\begin{exercise}[$\star\star\star$]\label{exo:dv:black-scholes-three-ways:7}
\emph{Programmation.} Avec \texttt{implied\_vol}, trouver la volatilité pour laquelle le
call de référence (spot 100, prix d'exercice 100, un an, $r=5\%$) vaut exactement 10.00.
\end{exercise}

\begin{exercise}[$\star\star\star$]\label{exo:dv:black-scholes-three-ways:8}
\emph{Repérer l'erreur.} «~Notre Monte Carlo à 10\,000 trajectoires donne 10.64 pour le
call de référence, soit 0.19 au-dessus de la formule. La formule ignore les queues épaisses~: nous
devrions coter 10.64.~»
\end{exercise}

\section{Problème~: Un prix, trois voies}

\begin{problem}[{Problème du week-end --- calculer trois fois le prix de référence du desk}]\label{pb:dv:black-scholes-three-ways:1}
Le call de référence de l'\cref{ex:dv:black-scholes-three-ways:ref} ($S=K=100$,
$T=1$, $r=5\%$, $q=0$, $\sigma=20\%$) doit être évalué par chacune des méthodes du
chapitre, et les réponses comparées.

\medskip\noindent\textbf{Partie I --- Réplication.}
\begin{enumerate}
\item Donner le delta, le gamma et le theta (par an) du call.
\item Vérifier numériquement que le call satisfait l'\omterm{def:dv:black-scholes-three-ways:pde}{équation de Black--Scholes}.
\item Le long de la trajectoire simulée du chapitre, quel est l'écart entre le
\omterm{def:dv:no-arbitrage-and-the-fundamental-theorems:claim}{portefeuille de réplication} et l'option à l'échéance~?
\item Pourquoi l'écart n'est-il pas nul, et pourquoi son signe n'a-t-il pas d'importance pour le prix~?
\item Où est passée la dérive de 12~\% de la trajectoire~?
\end{enumerate}

\medskip\noindent\textbf{Partie II --- La voie martingale.}
\begin{enumerate}[resume]
\item Donner $d_1$, $d_2$, $\Phi(d_1)$ et $\Phi(d_2)$.
\item Donner le call et le put par la formule, et vérifier la parité call-put.
\item Donner les probabilités d'exercice risque-neutre et réelle ($\mu=12\%$).
\item Donner le prix du call digital.
\item Comparer le call à l'approximation à la monnaie $0.4S\sigma\sqrt T$~; pourquoi
sous-estime-t-elle le prix ici~?
\end{enumerate}

\medskip\noindent\textbf{Partie III --- Méthodes numériques.}
\begin{enumerate}[resume]
\item Donner le prix de l'\omterm{def:dv:the-binomial-model:crr}{arbre de Cox--Ross--Rubinstein} à 1\,000 pas et son erreur.
\item Donner le prix Monte Carlo avec $10^6$ trajectoires, son erreur type et son erreur.
\item Combien de trajectoires faudrait-il pour atteindre la précision de l'arbre~?
\item Donner le prix Monte Carlo avec $10^4$ trajectoires et son erreur type.
\item Quelle méthode retiendrait-on pour 10\,000 calls européens, et laquelle pour un
put américain~?
\end{enumerate}

\medskip\noindent\textbf{Partie IV --- Jugement.}
\begin{enumerate}[resume]
\item Quelle hypothèse du modèle est la plus visiblement mise en défaut sur le marché~?
\item Pourquoi les desks cotent-ils encore en volatilité de Black--Scholes~?
\item Que changerait-il dans la formule si l'action versait un rendement de dividende de 2~\%~?
\item Énoncer le \emph{résultat nommé}~: le prix de référence par les trois voies, avec
l'erreur de chacune.
\item En une phrase~: pourquoi la dérive n'apparaît-elle pas~?
\end{enumerate}
\end{problem}

\section{Questions d'entretien}

\begin{interviewq}[$\star$ \iqroles{chercheur, trader}]\label{iq:dv:black-scholes-three-ways:1}
Dérivez l'\omterm{def:dv:black-scholes-three-ways:pde}{équation de Black--Scholes}.
\end{interviewq}

\begin{interviewq}[$\star$ \iqroles{trader}]\label{iq:dv:black-scholes-three-ways:2}
Un call à un an, à la monnaie, sur une action à 200 avec une volatilité de 30~\%~:
combien vaut-il à peu près~? (Ignorez les taux.)
\end{interviewq}

\begin{interviewq}[$\star\star$ \iqroles{chercheur, banque}]\label{iq:dv:black-scholes-three-ways:3}
Que sont $\Phi(d_1)$ et $\Phi(d_2)$~?
\end{interviewq}

\begin{interviewq}[$\star\star$ \iqroles{chercheur}]\label{iq:dv:black-scholes-three-ways:4}
Pourquoi le rendement espéré de l'action n'apparaît-il pas dans la formule~? Apparaîtrait-il
si vous ne pouviez pas négocier l'action~?
\end{interviewq}

\begin{interviewq}[$\star\star$ \iqroles{trader, risque}]\label{iq:dv:black-scholes-three-ways:5}
Énumérez les hypothèses du modèle. Quelle défaillance coûte le plus cher à un desk de volatilité~?
\end{interviewq}

\begin{interviewq}[$\star\star\star$ \iqroles{développeur}]\label{iq:dv:black-scholes-three-ways:6}
Vous devez calculer des volatilités implicites pour 50 millions de cotations d'options par jour,
y compris les options profondément hors de la monnaie et celles proches de l'échéance. Comment rendez-vous le
solveur robuste et rapide~?
\end{interviewq}
